\SourceLeaf{195}bringen könnte. Vielleicht werde ich, wenn ich einmal
mehr Muße habe als unter meinen gegenwärtigen Umständen, selbst
in Zukunft meine Untersuchungen bekannt machen.

Mit Hochachtung verbleibe ich
\begin{center}
\begin{tabular}{@{}l@{\qquad}r@{}}
 & Euer Wohlgeboren\\
Göttingen, den 8. November & ergebenster Diener\\
\multicolumn{1}{c}{1824.} & \textsc{C. F. Gauss.}
\end{tabular}
\end{center}
\Needspace{9\baselineskip}
\medskip

\begin{center}
\textsc{Gauss} an \textsc{Olbers}. Göttingen, 3. Mai 1827.
\end{center}

.\ .\ .\ .\ .\ .\ .\ Vor etwa 6 Wochen hatte ich das Vergnügen, Herrn \textsc{Dirichlet},
von dem ich Ihnen, glaube ich, schon einmal geschrieben habe, hier persönlich
kennenzulernen. Ich erwähnte ihm gegenüber den schlechten Aufsatz von
\textsc{Ivory}. Er kannte ihn selbst nicht, sagte mir aber, Herr \textsc{Fourier}
habe ihm gesagt, dass er »Unsinn« sei und dass er (F.) sehr über einen
lobhudelnden Artikel darüber im Férussac gelacht habe. Ungefähr ebenso wie
ich habe er auch über sein Mémoire zur Gleichgewichtsfigur einer
rotierenden homogenen Flüssigkeit geurteilt. Zu meiner Überraschung stellte ich auch in
mehreren anderen Fragen zu Personen und Sachen eine außerordentliche
Übereinstimmung zwischen Herrn \textsc{Fouriers} Urteil und meinem fest, z.\,B.
über imaginäre Größen, über die Unbeweisbarkeit der Geometrie a priori und so
weiter. Nachdem ich Herrn \textsc{Dirichlet} z.\,B. meine Ansicht dazu kurz
angedeutet hatte, sagte er, dass ihm \textsc{Fourier} seine Ansicht fast mit denselben
Worten gesagt habe.
.\ .\ .\ .\ .\ .

